\documentclass[aspectratio=169]{beamer}
\usetheme{Madrid}
\usecolortheme{default}
\usepackage[utf8]{inputenc}
\usepackage[T5]{fontenc}
\usepackage{graphicx}
\usepackage{hyperref}
\usepackage{booktabs}
\usepackage{tcolorbox}

\setbeamertemplate{caption}[numbered]
\renewcommand{\figurename}{Hình}
\renewcommand{\thefigure}{1.\arabic{figure}}

\title[Tổng quan Học Sâu]{Học Sâu (Deep Learning) \\ \vspace{0.5cm} \Large Cẩm nang Học tập: Từ Nền tảng đến AI Tạo sinh}
\author{Đại học Đông Á}
\date{\today}

\begin{document}

\begin{frame}
    \titlepage
\end{frame}

\begin{frame}{Nội dung chương}
    \tableofcontents[hideallsubsections]
\end{frame}

% ==============================
\section{Mở đầu: Bản đồ Tri thức Môn Học Sâu}
% ==============================

\begin{frame}{Bài toán sư phạm trong Học Sâu}
    \begin{itemize}
        \item \textbf{Thách thức cốt lõi:} Cuốn sách \textit{Deep Learning with Python} (2025) chứa một lượng kiến thức khổng lồ.
        \item \textbf{Nguy cơ:} Sinh viên dễ bị ngợp và mất phương hướng nếu chỉ tiếp cận tuần tự mà không có bức tranh tổng thể.
        \item \textbf{Mục tiêu:} Xây dựng một lộ trình tóm tắt giúp sinh viên:
        \begin{itemize}
            \item Không bị choáng ngợp trước khối lượng kiến thức.
            \item Thấy rõ sợi dây liên kết logic với các môn học tiền đề đã hoàn thành (AI, ML, CV).
        \end{itemize}
    \end{itemize}
\end{frame}

\begin{frame}{Hướng tiếp cận: Tree of Thought}
    \begin{itemize}
        \item Áp dụng mô hình \textbf{Tree of Thought} (Cây tư duy) để phân rã và thiết kế lộ trình.
        \item Lộ trình được chia thành \textbf{3 Phần chính}, kết nối hữu cơ với nhau.
    \end{itemize}
    \begin{center}
        \begin{tcolorbox}[width=0.8\textwidth, colback=blue!5, colframe=blue!40!black, title=Bản đồ Tri thức]
            \begin{itemize}
                \item \textbf{Phần 1:} Bản đồ Tri thức nền tảng
                \item \textbf{Phần 2:} Học sâu cơ bản \& Thị giác
                \item \textbf{Phần 3:} Chuỗi thời gian, NLP \& AI Tạo sinh tiên tiến
            \end{itemize}
        \end{tcolorbox}
    \end{center}
\end{frame}

\begin{frame}{Giới thiệu Phần 1: Bản đồ Tri thức nền tảng}
    \begin{itemize}
        \item \textbf{Nội dung:} Tóm tắt lại 3 trụ cột tri thức đã học.
        \begin{itemize}
            \item Trí tuệ Nhân tạo (AI)
            \item Học máy (ML)
            \item Thị giác Máy tính (CV)
        \end{itemize}
        \item \textbf{Mục tiêu:} Giúp sinh viên hệ thống hóa \textit{“những gì đã biết”} và thấy được nhu cầu tất yếu dẫn đến Học sâu.
    \end{itemize}
\end{frame}

\begin{frame}{Giới thiệu Phần 2: Hành trình Chinh phục Học sâu - Tập 1}
    \begin{itemize}
        \item \textbf{Nội dung:} Chinh phục Chương 1 đến Chương 12 của cuốn sách.
        \item \textbf{Trọng tâm:}
        \begin{itemize}
            \item Nền tảng mạng nơ-ron.
            \item Thư viện Keras / TensorFlow / PyTorch / JAX.
            \item Các bài toán cốt lõi của Thị giác Máy tính.
        \end{itemize}
        \item \textbf{Ứng dụng:} Phân loại ảnh, Phân đoạn ảnh, Phát hiện đối tượng.
    \end{itemize}
\end{frame}

\begin{frame}{Giới thiệu Phần 3: Hành trình Chinh phục Học sâu - Tập 2}
    \begin{itemize}
        \item \textbf{Nội dung:} Chinh phục Chương 13 đến Chương 20 của cuốn sách.
        \item \textbf{Trọng tâm:}
        \begin{itemize}
            \item Các mô hình tuần tự.
            \item Xử lý ngôn ngữ tự nhiên (NLP) với Transformers.
            \item Các kiến trúc AI tạo sinh (VAE, Diffusion, LLMs/GPT).
        \end{itemize}
        \item \textbf{Kỹ năng:} Thực chiến đưa mô hình ra thực tế.
    \end{itemize}
\end{frame}

% ==============================
\section{Phần 1: Bản đồ Tri thức nền tảng (Prior Knowledge)}
% ==============================

\begin{frame}
    \vfill
    \centering
    \begin{beamercolorbox}[sep=8pt,center,shadow=true,rounded=true]{title}
        \usebeamerfont{title}\Large PHẦN 1: BẢN ĐỒ TRI THỨC NỀN TẢNG
    \end{beamercolorbox}
    \vspace{0.5cm}
    \textit{Kích hoạt lại vùng bộ nhớ của ba môn học nền tảng đã tích lũy.}
    \vfill
\end{frame}

\begin{frame}{Trụ cột 1: Trí tuệ Nhân tạo (Artificial Intelligence)}
    \begin{itemize}
        \item \textbf{Tài liệu hệ thống:} \textit{Artificial Intelligence: A Modern Approach} (Stuart J. Russell, Peter Norvig).
        \item Mở ra thế giới quan tổng quan nhất về AI.
        \item \textbf{Khái niệm cốt lõi:} Tác tử hợp lý (Rational Agent) — Mô hình chuẩn (Standard Model).
        \item \textbf{Triết lý:} Xây dựng các hệ thống có khả năng nhận thức môi trường và đưa ra hành động tối ưu nhất.
    \end{itemize}
\end{frame}

\begin{frame}{Tác tử và Môi trường}
    \begin{itemize}
        \item Một thực thể nhận thông tin đầu vào (Percepts) qua \textbf{cảm biến} (Sensors).
        \item Tác động ngược lại môi trường qua \textbf{bộ chấp hành} (Actuators).
        \item \textbf{Sự hợp lý (Rationality):} Phụ thuộc vào hàm đo lường hiệu quả (Performance measure), tri thức nền tảng và chuỗi thông tin nhận được.
    \end{itemize}
\end{frame}

\begin{frame}{Giải quyết vấn đề bằng Tìm kiếm (Search)}
    \begin{itemize}
        \item \textbf{Tìm kiếm mù (Uninformed Search):} BFS, DFS, Uniform-cost.
        \item \textbf{Tìm kiếm kinh nghiệm (Heuristic search):} Sử dụng hàm Heuristic để hướng dẫn quá trình tìm kiếm.
        \item \textbf{Ví dụ nổi bật:} Giải thuật $A^*$ trong tìm kiếm không gian trạng thái phức tạp.
    \end{itemize}
\end{frame}

\begin{frame}{Ra quyết định và Học tăng cường cơ bản}
    \begin{itemize}
        \item \textbf{Ra quyết định bất định:}
        \begin{itemize}
            \item Thuyết xác suất và Thuyết hữu dụng (Utility theory).
            \item Tối đa hóa độ hữu dụng kỳ vọng. Tiền đề cho Mạng Bayes.
        \end{itemize}
        \item \textbf{Học tăng cường (RL):}
        \begin{itemize}
            \item Quá trình quyết định Markov (MDPs).
            \item Học qua cơ chế thưởng/phạt từ môi trường để tìm \textbf{Chính sách (Policy)} tối ưu.
        \end{itemize}
        \item \textbf{Sự dịch chuyển lịch sử:} Từ hệ chuyên gia (expert systems) sang tác tử tự học hỏi.
    \end{itemize}
\end{frame}

\begin{frame}{Trụ cột 2: Học máy (Machine Learning)}
    \begin{itemize}
        \item \textbf{Tài liệu hệ thống:} \textit{Hands-On Machine Learning with Scikit-Learn and PyTorch} (Aurélien Géron).
        \item AI cho tư duy hệ thống, ML cung cấp bộ công cụ toán học và lập trình thực chiến.
        \item Máy tính tự học các quy luật từ dữ liệu mà không cần lập trình tường minh.
        \item \textbf{Hai khối cốt lõi:} ML cổ điển (Scikit-Learn) và Chuyển tiếp (PyTorch).
    \end{itemize}
    \begin{figure}
        \centering
        \includegraphics[width=0.5\textwidth,keepaspectratio]{example-image} % TODO: Thay ảnh ML bằng đường dẫn: D:\DongAUniversity\TÀI LIỆU DẠY HỌC_2024-2025\Môn Máy học_2026\machineLearningWeb\...
        \caption{Quy trình Học máy cơ bản}
    \end{figure}
\end{frame}

\begin{frame}{Quy trình chuẩn và Học có giám sát (Supervised)}
    \begin{itemize}
        \item \textbf{Quy trình chuẩn:}
        \begin{itemize}
            \item Khảo sát, làm sạch (Cleaning).
            \item Chuẩn hóa tính năng (Feature Scaling).
            \item Đường ống xử lý (Pipelines).
        \end{itemize}
        \item \textbf{Các bài toán cốt lõi:}
        \begin{itemize}
            \item \textbf{Hồi quy (Regression):} Dự đoán giá trị liên tục (Linear Regression, Random Forests).
            \item \textbf{Phân loại (Classification):} Dự đoán nhãn (SVM, SGDClassifier).
        \end{itemize}
    \end{itemize}
\end{frame}

\begin{frame}{Đánh giá và Học không giám sát}
    \begin{itemize}
        \item \textbf{Đánh giá mô hình phân loại:} Ma trận nhầm lẫn (Confusion Matrix), Precision, Recall.
        \item \textbf{Học không giám sát (Unsupervised):}
        \begin{itemize}
            \item \textbf{Phân cụm (Clustering):} K-Means, DBSCAN.
            \item \textbf{Giảm chiều dữ liệu (Dimensionality Reduction):} PCA giải quyết ``lời nguyền chiều dữ liệu''.
        \end{itemize}
    \end{itemize}
\end{frame}

\begin{frame}{Tinh chỉnh và Kiểm soát mô hình}
    \begin{itemize}
        \item \textbf{Tinh chỉnh:}
        \begin{itemize}
            \item Đánh giá chéo (Cross-validation).
            \item Tìm kiếm siêu tham số: Grid Search, Randomized Search.
        \end{itemize}
        \item \textbf{Bẫy Overfitting \& Underfitting:}
        \begin{itemize}
            \item Sự đánh đổi Bias/Variance (Độ chệch \& Phương sai).
            \item Sử dụng Regularization (L1/L2) để kiểm soát sức mạnh mô hình.
        \end{itemize}
    \end{itemize}
\end{frame}

\begin{frame}{Khối chuyển tiếp sang Mạng nơ-ron (PyTorch)}
    \begin{itemize}
        \item \textbf{Perceptron và MLP:} Đi từ cấu trúc nơ-ron sinh học sang nơ-ron nhân tạo và mạng MLP (Multilayer Perceptron).
        \item \textbf{Huấn luyện mạng sâu:}
        \begin{itemize}
            \item Lan truyền ngược (Backpropagation).
            \item Hiện tượng triệt tiêu/bùng nổ gradient.
            \item Các bộ khởi tạo trọng số (Glorot/He).
        \end{itemize}
    \end{itemize}
\end{frame}

\begin{frame}{Trụ cột 3: Thị giác Máy tính (Computer Vision)}
    \begin{itemize}
        \item Giúp máy tính ``nhìn'' và ``hiểu'' thông qua camera.
        \item Lộ trình chuyển hóa:
        \begin{center}
            \textbf{Ảnh thô $\rightarrow$ Tiền xử lý $\rightarrow$ Trích chọn đặc trưng $\rightarrow$ Nhận dạng}
        \end{center}
        \item \textbf{Biểu diễn ảnh số:} Ảnh là mảng đa chiều (Arrays), mô hình camera vật lý.
    \end{itemize}
    \begin{figure}
        \centering
        \includegraphics[width=0.5\textwidth,keepaspectratio]{example-image} % TODO: Thay ảnh CV bằng đường dẫn: D:\DongAUniversity\TÀI LIỆU DẠY HỌC_2024-2025\Trí tuệ nhân tạo_UDA_2025\webTTNT_2026\...
        \caption{Ví dụ về xử lý Thị giác máy tính}
    \end{figure}
\end{frame}

\begin{frame}{Tiền xử lý và Phân đoạn ảnh}
    \begin{itemize}
        \item \textbf{Tiền xử lý \& Lọc nhiễu:}
        \begin{itemize}
            \item Phép toán điểm ảnh, Cân bằng Histogram.
            \item Bộ lọc Gaussian (Làm mờ), Sobel, Laplace (Phát hiện biên).
        \end{itemize}
        \item \textbf{Phân đoạn ảnh:}
        \begin{itemize}
            \item Phân tách nền bằng ngưỡng Otsu.
            \item Thuật toán Watershed (phân tách vật thể chồng lấn).
            \item Thuật toán tương tác GrabCut.
        \end{itemize}
    \end{itemize}
\end{frame}

\begin{frame}{Phát hiện đặc trưng và Phân lớp cổ điển}
    \begin{itemize}
        \item \textbf{Phát hiện \& Mô tả đặc trưng:}
        \begin{itemize}
            \item Trích xuất điểm bất biến: Harris Corner, SIFT, SURF, ORB.
            \item Mô tả cấu trúc: HOG, LBP.
            \item So khớp (Feature Matching): KNN, FLANN.
        \end{itemize}
        \item \textbf{Phân lớp ảnh cổ điển:}
        \begin{itemize}
            \item Kết hợp Bag-of-Visual-Words với các bộ phân lớp mạnh (SVM, KNN) để nhận diện vật thể.
        \end{itemize}
    \end{itemize}
\end{frame}

\begin{frame}{Chuyển dịch sang Deep Learning trong CV}
    \begin{itemize}
        \item Mạng tích chập CNN (Convolutional Neural Networks).
        \item Kỹ thuật \textbf{Transfer Learning} với các mạng nổi tiếng (AlexNet, VGG16, ResNet).
        \item \textbf{Phát hiện \& Đánh giá chuyển động:}
        \begin{itemize}
            \item Vẽ khung bao (Bounding Box) thời gian thực: YOLO, SSD, Faster R-CNN.
            \item Phép biến đổi Homography.
            \item Tính toán Luồng quang (Optical Flow) bằng thuật toán Lucas-Kanade.
        \end{itemize}
    \end{itemize}
\end{frame}

\begin{frame}{Sợi dây liên kết: Tại sao cần Học sâu?}
    \begin{itemize}
        \item \textbf{Câu hỏi:} Nếu đã có SIFT, HOG (ảnh), mạng Bayes (văn bản)... tại sao lại cần DL?
        \item \textbf{Bản chất vấn đề cũ:} Con người phải tự thiết kế bộ trích chọn đặc trưng $\rightarrow$ Tốn công, không tối ưu cho dữ liệu khổng lồ.
        \item \textbf{Giải pháp Học sâu:} \textbf{Học biểu diễn (Representation Learning)}.
        \begin{itemize}
            \item Mô hình \textit{tự động} bóc tách các lớp biểu diễn từ thô đến tinh.
            \item Sử dụng các bộ lọc tham số được tối ưu bằng toán học (Lan truyền ngược).
        \end{itemize}
    \end{itemize}
\end{frame}

% ==============================
\section{Phần 2: Hành trình Chinh phục Học sâu (Chương 1-12)}
% ==============================

\begin{frame}
    \vfill
    \centering
    \begin{beamercolorbox}[sep=8pt,center,shadow=true,rounded=true]{title}
        \usebeamerfont{title}\Large PHẦN 2: HÀNH TRÌNH CHINH PHỤC HỌC SÂU
    \end{beamercolorbox}
    \vspace{0.5cm}
    \textit{Chương 1 - 12: Từ cơ bản Keras đến các mạng tích chập CV tiên tiến.}
    \vfill
\end{frame}

\begin{frame}{Phần 2: Mục tiêu và Lộ trình}
    \begin{itemize}
        \item \textbf{Mục tiêu:} Chuyển hóa tư duy từ Học máy cổ điển sang Học sâu thông qua thư viện Keras 3.
        \item \textbf{Cấu trúc (4 Cột mốc):}
        \begin{enumerate}
            \item Từ Toán học đến Code (Ch 1-3)
            \item Lập trình dự án \& Kiểm soát Overfitting (Ch 4-7)
            \item Vén màn hộp đen của CNNs (Ch 8-10)
            \item Thế giới Thị giác nâng cao (Ch 11-12)
        \end{enumerate}
    \end{itemize}
\end{frame}

\begin{frame}{Cột mốc 1: Xây dựng nền móng (Ch 1-3)}
    \begin{itemize}
        \item \textbf{Chương 1: Học sâu là gì?}
        \begin{itemize}
            \item \textit{Sự dịch chuyển biểu diễn:} Lập trình cổ điển (nhập luật) vs ML (tự tìm luật) vs DL (tự học biểu diễn đa tầng).
            \item Cảnh báo bong bóng truyền thông, khẳng định giá trị thực của mạng nơ-ron.
        \end{itemize}
        \item \textbf{Chương 2: Các khối toán học}
        \begin{itemize}
            \item \textbf{Tensors:} Scalars, Vectors, Matrices.
            \item \textbf{Tensor Operations:} Cộng, Broadcasting, Dot product, Reshaping.
            \item \textbf{Động cơ Tối ưu hóa:} Gradient, SGD, Backpropagation (Tự code vòng lặp từ đầu).
        \end{itemize}
    \end{itemize}
\end{frame}

\begin{frame}{Hệ sinh thái Frameworks (Chương 3)}
    \begin{itemize}
        \item \textbf{TensorFlow, PyTorch, JAX:}
        \begin{itemize}
            \item TensorFlow: Mạnh về deploy công nghiệp.
            \item PyTorch: Tiêu chuẩn học thuật (Hugging Face).
            \item JAX: Hiệu năng cao cho tính toán khoa học.
        \end{itemize}
        \item \textbf{Keras 3 - Kẻ thống nhất:}
        \begin{itemize}
            \item Giao diện API cấp cao chạy mượt trên cả 3 backend.
            \item Khái niệm: Lớp (Layers), Mô hình, Compile (Loss, Optimizer), fit().
        \end{itemize}
    \end{itemize}
\end{frame}

\begin{frame}{Cột mốc 2: Thực chiến Dự án \& Quy trình chuẩn (Ch 4-7)}
    \begin{itemize}
        \item \textbf{Chương 4: Phân lớp \& Hồi quy (Dữ liệu Vector)}
        \begin{itemize}
            \item \textit{Phân loại nhị phân (IMDb):} Cấu hình ngõ ra Sigmoid + Binary Crossentropy.
            \item \textit{Phân loại đa lớp (Reuters):} Ngõ ra Softmax + Categorical Crossentropy.
            \item \textit{Hồi quy (California Housing):} Không có hàm kích hoạt cuối, dùng MSE \& K-fold.
        \end{itemize}
    \end{itemize}
\end{frame}

\begin{frame}{Chinh phục Tổng quát hóa (Chương 5)}
    \begin{itemize}
        \item \textbf{Optimization vs Generalization:}
        \begin{itemize}
            \item Độ chính xác 100\% trên tập train thường là dấu hiệu của Overfitting thảm họa.
        \end{itemize}
        \item \textbf{Đánh giá mô hình:} Tách biệt Training, Validation, Test. Mốc baseline.
        \item \textbf{Vũ khí chống Overfitting:}
        \begin{itemize}
            \item Regularization (L1/L2)
            \item Early stopping
            \item \textbf{Dropout} (Xóa ngẫu nhiên các nơ-ron).
        \end{itemize}
    \end{itemize}
\end{frame}

\begin{frame}{Quy trình Học máy \& Keras chuyên sâu (Ch 6-7)}
    \begin{itemize}
        \item \textbf{Chương 6: Universal Workflow}
        \begin{itemize}
            \item 1. Định nghĩa bài toán \& thước đo.
            \item 2. Phát triển mô hình: Xây baseline $\rightarrow$ Overfit cố ý $\rightarrow$ Tối ưu hóa (Regularization).
            \item 3. Triển khai (Deploy): Đóng gói, đẩy lên server/mobile, giám sát data drift.
        \end{itemize}
        \item \textbf{Chương 7: Đi sâu vào Keras}
        \begin{itemize}
            \item 3 API: Sequential, Functional (Mạng đồ thị phức tạp), Subclassing.
            \item Callbacks, TensorBoard, Custom Training Loops.
        \end{itemize}
    \end{itemize}
\end{frame}

\begin{frame}{Cột mốc 3: Kỷ nguyên CV - Sự tích chập (Ch 8-10)}
    \begin{center}
        \textbf{Ảnh thô $\rightarrow$ Tích chập/Pooling $\rightarrow$ Đặc trưng hình học $\rightarrow$ Phân lớp}
    \end{center}
    \begin{itemize}
        \item \textbf{Chương 8: Giới thiệu ConvNets}
        \begin{itemize}
            \item \textit{Tích chập (Convolution):} Tự động học cửa sổ trượt (kernels) phát hiện đặc trưng (tính bất biến dịch chuyển).
            \item \textit{Max-pooling:} Giảm chiều, tăng trường thụ cảm (receptive field).
            \item \textit{Dữ liệu nhỏ:} Data Augmentation, Feature Extraction, Fine-tuning.
        \end{itemize}
    \end{itemize}
\end{frame}

\begin{frame}{Kiến trúc \& Trực quan hóa ConvNet (Ch 9-10)}
    \begin{itemize}
        \item \textbf{Chương 9: Mẫu thiết kế (Patterns)}
        \begin{itemize}
            \item \textit{Residual connections:} Tạo đường tắt tránh triệt tiêu gradient.
            \item \textit{Batch normalization:} Ổn định và tăng tốc độ hội tụ.
            \item \textit{Depthwise separable convolutions:} Tách tích chập không gian \& kênh, giảm tham số (MobileNet).
        \end{itemize}
        \item \textbf{Chương 10: Xóa bỏ ``hộp đen''}
        \begin{itemize}
            \item Trực quan hóa các kích hoạt trung gian (Intermediate activations).
            \item Trực quan hóa bộ lọc (Filters visualization) bằng Gradient Ascent.
            \item \textbf{Grad-CAM:} Tạo biểu đồ nhiệt chỉ ra vùng pixel quyết định kết quả.
        \end{itemize}
    \end{itemize}
\end{frame}

\begin{frame}{Cột mốc 4: Thị giác nâng cao (Ch 11-12)}
    \begin{itemize}
        \item \textbf{Chương 11: Phân đoạn ảnh (Image Segmentation)}
        \begin{itemize}
            \item Bóc tách ranh giới pixel của từng vật thể.
            \item \textit{Encoder-Decoder:} Nén ngữ cảnh $\rightarrow$ Giải nén (Conv2DTranspose) gán nhãn pixel.
            \item \textit{Segment Anything Model (SAM):} Mô hình nền tảng, điều khiển bằng prompt.
        \end{itemize}
        \item \textbf{Chương 12: Phát hiện đối tượng (Object Detection)}
        \begin{itemize}
            \item Vẽ khung bao (Bounding boxes).
            \item \textit{Single-stage (YOLO, SSD)} vs \textit{Two-stage (R-CNN)}.
            \item Thực hành tự xây dựng YOLO, dự đoán xác suất qua các grid cells.
        \end{itemize}
    \end{itemize}
\end{frame}

\begin{frame}{Chuyển tiếp sang Phần 3}
    \begin{itemize}
        \item Cuối Chương 12: Làm chủ \textbf{Dữ liệu Vector} (Ch 4-7) \& \textbf{Ảnh tĩnh} (Ch 8-12).
        \item Nhưng thế giới thực vận hành theo \textbf{dòng chảy thời gian} (tài chính, IoT) và \textbf{ngôn ngữ tự nhiên}.
        \item Đã đến lúc chuyển sang Dữ liệu tuần tự (RNN, Transformers) và AI Tạo sinh.
    \end{itemize}
\end{frame}

% ==============================
\section{Phần 3: Dòng chảy Thời gian, NLP \& AI Tạo sinh}
% ==============================

\begin{frame}
    \vfill
    \centering
    \begin{beamercolorbox}[sep=8pt,center,shadow=true,rounded=true]{title}
        \usebeamerfont{title}\Large PHẦN 3: DÒNG CHẢY THỜI GIAN, NLP \& AI TẠO SINH
    \end{beamercolorbox}
    \vspace{0.5cm}
    \textit{Chương 13 - 20: LSTM, Attention, Transformers, LLMs và Diffusion.}
    \vfill
\end{frame}

\begin{frame}{Phần 3: Mục tiêu và Lộ trình}
    \begin{itemize}
        \item \textbf{Mục tiêu:} Tiếp cận những kiến trúc hiện đại nhất định hình cách mạng AI.
        \item \textbf{Cấu trúc (4 Cột mốc):}
        \begin{enumerate}
            \setcounter{enumi}{4}
            \item Dòng chảy Thời gian \& Chữ viết (Ch 13-14)
            \item Kỷ nguyên Transformers \& GPT (Ch 15-16)
            \item Thế giới AI Tạo ảnh (Ch 17)
            \item Đưa AI ra Thực tế \& Tương lai (Ch 18-20)
        \end{enumerate}
    \end{itemize}
\end{frame}

\begin{frame}{Cột mốc 5: Dòng chảy Thời gian \& Chữ viết (Ch 13-14)}
    \begin{itemize}
        \item \textbf{Chương 13: Dự báo chuỗi thời gian}
        \begin{itemize}
            \item Tính tuần tự và phụ thuộc thời gian (ví dụ: thời tiết).
            \item \textbf{Mạng RNNs:} Đặc biệt là \textbf{LSTM} và \textbf{GRU} — cấu trúc cổng kiểm soát giúp lưu giữ thông tin dài hạn.
            \item Dropout chuyên dụng cho dữ liệu chuỗi.
        \end{itemize}
        \item \textbf{Chương 14: Phân lớp văn bản}
        \begin{itemize}
            \item \textbf{Tokenization:} Biến ngôn ngữ thành số học.
            \item Bag-of-words (thống kê) vs Sequence models (ngữ cảnh).
            \item \textbf{Word Embeddings (Nhúng từ):} Không gian hình học của ngữ nghĩa (khoảng cách vector tương đồng nghĩa).
        \end{itemize}
    \end{itemize}
\end{frame}

\begin{frame}{Cột mốc 6: Kỷ nguyên Transformers \& GPT (Ch 15-16)}
    \begin{itemize}
        \item \textbf{Chương 15: Mô hình ngôn ngữ \& Transformer}
        \begin{itemize}
            \item Bản chất mô hình ngôn ngữ: Dự đoán từ tiếp theo.
            \item \textbf{Cơ chế Self-Attention (Tự chú ý):} Nhìn vào toàn bộ câu để bắt ngữ cảnh cùng một lúc (thay thế RNNs).
            \item Kiến trúc Transformer: Encoder, Decoder, Positional embedding.
        \end{itemize}
        \item \textbf{Chương 16: Sinh văn bản (Text Generation)}
        \begin{itemize}
            \item Huấn luyện mini-GPT từ đầu (Decoder-only).
            \item \textbf{Chiến lược lấy mẫu:} Greedy, Random, Top-K (Tạo sinh tự nhiên, không lặp).
            \item Tinh chỉnh chỉ dẫn (Instruction fine-tuning) \& RLHF.
        \end{itemize}
    \end{itemize}
\end{frame}

\begin{frame}{Cột mốc 7: Thế giới Mê hoặc của AI Tạo ảnh (Ch 17)}
    \begin{itemize}
        \item \textbf{Không gian tiềm ẩn (Latent spaces):}
        \begin{itemize}
            \item Nén ảnh thành vector, biến đổi tọa độ để thay đổi thuộc tính (thêm kính, đổi nụ cười).
        \end{itemize}
        \item \textbf{VAEs (Variational Autoencoders):}
        \begin{itemize}
            \item Mã hóa thành phân phối xác suất, giải mã sinh ảnh mới.
        \end{itemize}
        \item \textbf{Diffusion models (Mô hình Khuếch tán):}
        \begin{itemize}
            \item Phép thuật của Midjourney, Stable Diffusion.
            \item Học cách loại bỏ nhiễu (denoising) từng bước từ nhiễu hạt hoàn toàn dựa trên Text-to-image.
        \end{itemize}
    \end{itemize}
\end{frame}

\begin{frame}{Cột mốc 8: Nâng cấp Chuẩn Sản xuất (Ch 18-20)}
    \begin{itemize}
        \item \textbf{Chương 18: Kỹ thuật thực chiến}
        \begin{itemize}
            \item Tối ưu hóa siêu tham số bằng \textbf{KerasTuner}.
            \item Kết hợp mô hình (Ensembling), Huấn luyện phân tán, Mixed-precision.
            \item Lượng tử hóa (Quantization) cho thiết bị yếu.
        \end{itemize}
        \item \textbf{Chương 19: Tương lai của AI}
        \begin{itemize}
            \item Hạn chế: Mạng nơ-ron thực chất là Nội suy (Interpolation), yếu về suy luận logic (Extrapolation).
            \item Hướng mới: ARC-AGI, Tổng hợp chương trình (Program Synthesis).
        \end{itemize}
        \item \textbf{Chương 20: Tổng kết hành trình}
        \begin{itemize}
            \item Liên kết hệ thống: Dense $\leftrightarrow$ ConvNets $\leftrightarrow$ RNNs $\leftrightarrow$ Transformers.
            \item Học suốt đời: arXiv, Kaggle.
        \end{itemize}
    \end{itemize}
\end{frame}

\begin{frame}{Kết luận: Khởi hành}
    \begin{itemize}
        \item \textbf{1. Niềm tin nền tảng:} Các em đã nắm 70\% gốc rễ từ AI, ML, CV.
        \item \textbf{2. Hành động thực chiến:} Làm chủ Keras 3 \& CV sâu sắc.
        \item \textbf{3. Mở rộng tầm nhìn:} Đỉnh cao công nghệ (LLMs, Diffusion).
        \item Sự nghiệp kỹ sư AI bắt đầu từ hôm nay. Chúc các em học tập hiệu quả!
    \end{itemize}
\end{frame}

\end{document}
